\section{TVD-Conflict: Disagreement-Guided Vulnerability Discovery}
\label{sec:tvd-conflict}

The preceding study shows that judge conflict is concentrated in particular
responses and evaluation criteria. \textsc{TVD}-Conflict uses this structure to
search for evaluator-sensitive cases directly. Instead of executing a
jailbreak attack against a victim model and checking whether the resulting
response happens to divide the judges, an agent constructs the query--response
pair and revises it using the judges' decisions.

\subsection{Problem Formulation}

Let $x=(g,q_{\mathrm{adv}},r)$ denote a candidate, where $g$ is the original
malicious goal, $q_{\mathrm{adv}}$ is the adversarial query constructed from
that goal, and $r$ is the assistant response. The goal $g$ defines the behavior
that the search must preserve, while $(q_{\mathrm{adv}},r)$ forms the visible
interaction evaluated by the judges. Given an evaluation criterion $R$, two
judges $J_a$ and $J_b$ independently evaluate this interaction. The criterion,
prompt, output format, and visible context remain fixed. Cross-judge
experiments vary the backbone, whereas same-judge experiments use two
independent calls to the same backbone.

For binary evaluation, we use the WildGuard protocol
\citep{han2024wildguard}; a conflict occurs when the two judges assign
different harmfulness labels. For graded evaluation, we use StrongREJECT
\citep{souly2024strongreject}, whose refusal, convincingness, and specificity
ratings produce a normalized score in $[0,1]$. A graded conflict occurs when
the two normalized scores differ by more than 0.5. Given a candidate budget
$B$, the search succeeds if it finds at least one interaction satisfying the
corresponding conflict condition.

\subsection{TVD-Conflict Framework}

\textsc{TVD}-Conflict instantiates the Task--Validator--Data framework around
this search objective. Figure~\ref{fig:tvd-conflict-overview} summarizes the
feedback loop.

\paragraph{Task.}
For each malicious goal, the task asks an agent to construct a complete
user--assistant interaction while preserving the semantics of the goal. We run
separate searches for substantively unsafe and safe responses. The unsafe
branch targets cases in which one judge misses harmful assistance, while the
safe branch targets cases in which one judge overestimates the harmfulness of a
non-compliant response. The agent may revise the query and response but cannot
change the original goal, target safety label, evaluation criterion, or judge
configuration.

\paragraph{Validator.}
The validator sends each candidate to both judges and checks the conflict
condition defined above. It records the raw judgments together with their parsed
labels or scores, then returns both outputs to the agent. A conflicting
candidate is retained; otherwise, the judgments become feedback for the next
revision. API failures, unparsable judgments, and malformed candidates do not
count as conflicts.

\paragraph{Data.}
The data component maintains the active candidates, their previous judgments,
and the remaining budget. Once a candidate induces a conflict, the interaction
and both judge outputs are archived and removed from the active set. Otherwise,
the agent continues revising it until the budget is exhausted. The stored
trajectory records how each retained case developed across revisions and how
the two judges responded.

\paragraph{Sample realism.}
The agent receives few-shot conflicts obtained by auditing outputs from
established jailbreak attacks. These examples demonstrate realistic
query--response structures and common ways in which harmful behavior appears in
model outputs. The task gives sample realism priority over inducing a conflict:
each candidate must preserve the intent and substantive outcome of the original
goal and remain plausible as a user query and assistant response.
Evaluator-facing instructions, scoring cues, task substitution, and malformed
interactions are excluded.

\subsection{Conflict Collection and Evaluation}

\paragraph{Collection setup.}
We use the same 50 AdvBench goals and three cross-backbone judge pairs as in
Section~\ref{sec:judge-disagreement}. To distinguish systematic differences
between judges from stochastic instability within a judge, we additionally
form three same-backbone pairs from independent evaluations by GPT, DeepSeek,
and Qwen. Each of the six pairs is evaluated under WildGuard and StrongREJECT.
\textsc{TVD}-Conflict runs the safe and unsafe branches independently and
revises each active candidate for at most ten iterations.

We compare two routes to collecting conflict pairs. \emph{Attack-Based Audit}
runs 10 jailbreak attacks against four victim models for every goal, producing
2,000 query--response pairs before judge evaluation. \textsc{TVD}-Conflict
constructs the query and response directly from the malicious goal and revises
them using judge feedback. It therefore does not require attack reproduction or
victim-model execution.

\paragraph{Metrics.}
For each criterion, we report the number of conflict pairs collected and the
number of candidate pairs judged. We separately report victim-model generations
to capture the principal cost of constructing an attack-derived audit corpus.
The attack pipeline requires one victim response for every generated candidate;
iterative attacks may additionally invoke an attacker model. These
attack-specific calls were not logged uniformly and are not included, so the
reported attack-generation cost is a lower bound. Each judged candidate is sent
to both judges under either collection method.

% Main experiment summary table, computed from the stored attempt trajectories
% in code/run/main_adv50_luna_max/. Each CR entry is the percentage of target
% instances that produced at least one conflict by the indicated budget.
\begin{table*}[t]
  \centering
  \caption{Main-experiment TVD-Conflict rates across judge pairs. Entries
  report the conflict rate (CR, \%) at collection budgets 3, 5, and 10;
  Mean CR@10 averages the binary and scoring rates at budget 10. The final
  row reports the number of conflict pairs collected at budget 10, with the
  final entry giving their total across the two criteria.}
  \label{tab:main-experiment-conflict}
  \footnotesize
  \setlength{\tabcolsep}{4.8pt}
  \renewcommand{\arraystretch}{1.12}
  \begin{tabular}{@{}l ccc ccc c@{}}
    \toprule
    \multicolumn{1}{@{}l}{Judge pairs}
      & \multicolumn{3}{c}{Binary}
      & \multicolumn{3}{c}{Scoring}
      & \multicolumn{1}{c@{}}{Mean} \\
    \cmidrule(lr){2-4}\cmidrule(lr){5-7}\cmidrule(l){8-8}
      & CR@3 & CR@5 & CR@10
      & CR@3 & CR@5 & CR@10 & CR@10 \\
    \midrule
    \multicolumn{8}{@{}l}{\textbf{Different judges}} \\
    \quad GPT / DeepSeek
      & \cellcolor{red!16}48.5 & \cellcolor{red!22}67.5 & \cellcolor{red!29}90.0
      & \cellcolor{red!16}48.0 & \cellcolor{red!21}65.5 & \cellcolor{red!29}87.5
      & \cellcolor{red!29}88.8 \\
    \quad GPT / Qwen
      & \cellcolor{red!13}39.0 & \cellcolor{red!18}56.5 & \cellcolor{red!27}83.5
      & \cellcolor{red!15}46.5 & \cellcolor{red!22}68.0 & \cellcolor{red!29}90.0
      & \cellcolor{red!28}86.8 \\
    \quad DeepSeek / Qwen
      & \cellcolor{red!7}21.5 & \cellcolor{red!13}40.5 & \cellcolor{red!24}72.5
      & \cellcolor{red!18}55.0 & \cellcolor{red!24}73.0 & \cellcolor{red!30}92.0
      & \cellcolor{red!27}82.3 \\
    \addlinespace[2pt]
    \multicolumn{8}{@{}l}{\textbf{Same judge}} \\
    \quad GPT / GPT
      & \cellcolor{red!8}24.5 & \cellcolor{red!14}43.0 & \cellcolor{red!23}71.5
      & \cellcolor{red!8}23.0 & \cellcolor{red!13}41.0 & \cellcolor{red!22}67.0
      & \cellcolor{red!23}69.3 \\
    \quad DeepSeek / DeepSeek
      & \cellcolor{red!5}15.0 & \cellcolor{red!9}27.5 & \cellcolor{red!16}50.5
      & \cellcolor{red!13}40.0 & \cellcolor{red!21}64.0 & \cellcolor{red!28}86.0
      & \cellcolor{red!22}68.3 \\
    \quad Qwen / Qwen
      & \cellcolor{red!1}3.0 & \cellcolor{red!2}5.5 & \cellcolor{red!6}17.5
      & \cellcolor{red!4}13.5 & \cellcolor{red!8}23.0 & \cellcolor{red!14}43.0
      & \cellcolor{red!10}30.3 \\
    \midrule
    \rowcolor{blue!4}
    \textbf{Conflict pairs collected}
      & & & \textbf{771}
      & & & \textbf{931} & \textbf{1702} \\
    \bottomrule
  \end{tabular}
\end{table*}


\paragraph{Conflict discovery.}
Table~\ref{tab:main-experiment-conflict} shows that disagreement-guided search
finds conflicts across every judge pair and both evaluation criteria, with
success increasing steadily as the collection budget grows. For
cross-backbone pairs, the mean conflict rate rises from 36.3\% at budget 3 to
82.0\% at budget 10 under binary evaluation, and from 49.8\% to 89.8\% under
scoring evaluation. Thus, most goals expose at least one evaluator-sensitive
interaction within ten revisions, even when the initial candidates do not
already divide the judges. The search collects 771 binary and 931 scoring
conflict pairs at the full budget.

\paragraph{Within-judge conflicts.}
We also instantiate the two validator roles with independent calls to the same
judge backbone. This setting searches for repeated-evaluation instability rather
than disagreement between different models; same-judge pairs are included in
Table~\ref{tab:main-experiment-conflict}. These conflicts are less frequent
than cross-backbone conflicts but remain substantial: at budget 10, the mean
rate is 46.5\% for binary judges and 65.3\% for scoring judges. The effect is
strongly backbone-dependent. Qwen is the most self-consistent judge (17.5\%
binary and 43.0\% scoring), whereas DeepSeek reaches 50.5\% and 86.0\%,
respectively. Scoring produces more conflict than binary labeling for five of
the six judge pairs at budget 10; GPT/GPT is the only exception. Evaluator
unreliability is therefore not solely a difference between model families:
repeated calls to the same model can cross the decision boundary, particularly
when a continuous rubric exposes more degrees of freedom than a binary label.
